\documentclass{amsart}
% For dealing with references we use the comment environment
\usepackage{verbatim}
\newenvironment{reference}{\comment}{\endcomment}
%\newenvironment{reference}{}{}
\newenvironment{slogan}{\comment}{\endcomment}
\newenvironment{history}{\comment}{\endcomment}

% For commutative diagrams we use Xy-pic
\usepackage[all]{xy}

% We use 2cell for 2-commutative diagrams.
\xyoption{2cell}
\UseAllTwocells

% We use multicol for the list of chapters between chapters
\usepackage{multicol}

% This is generally recommended for better output
\usepackage{lmodern}
\usepackage[T1]{fontenc}

% For cross-file-references

% Package for hypertext links:
\usepackage{hyperref}

% For any local file, say "hello.tex" you want to link to please

% Theorem environments.
%
\theoremstyle{plain}
\newtheorem{theorem}[subsection]{Theorem}
\newtheorem{proposition}[subsection]{Proposition}
\newtheorem{lemma}[subsection]{Lemma}

\theoremstyle{definition}
\newtheorem{definition}[subsection]{Definition}
\newtheorem{example}[subsection]{Example}
\newtheorem{exercise}[subsection]{Exercise}
\newtheorem{situation}[subsection]{Situation}

\theoremstyle{remark}
\newtheorem{remark}[subsection]{Remark}
\newtheorem{remarks}[subsection]{Remarks}

\numberwithin{equation}{subsection}

% Macros
%
\def\lim{\mathop{\mathrm{lim}}\nolimits}
\def\colim{\mathop{\mathrm{colim}}\nolimits}
\def\Spec{\mathop{\mathrm{Spec}}}
\def\Hom{\mathop{\mathrm{Hom}}\nolimits}
\def\Ext{\mathop{\mathrm{Ext}}\nolimits}
\def\SheafHom{\mathop{\mathcal{H}\!\mathit{om}}\nolimits}
\def\SheafExt{\mathop{\mathcal{E}\!\mathit{xt}}\nolimits}
\def\Sch{\mathit{Sch}}
\def\Mor{\mathop{\mathrm{Mor}}\nolimits}
\def\Ob{\mathop{\mathrm{Ob}}\nolimits}
\def\Sh{\mathop{\mathit{Sh}}\nolimits}
\def\NL{\mathop{N\!L}\nolimits}
\def\CH{\mathop{\mathrm{CH}}\nolimits}
\def\proetale{{pro\text{-}\acute{e}tale}}
\def\etale{{\acute{e}tale}}
\def\QCoh{\mathit{QCoh}}
\def\Ker{\mathop{\mathrm{Ker}}}
\def\Im{\mathop{\mathrm{Im}}}
\def\Coker{\mathop{\mathrm{Coker}}}
\def\Coim{\mathop{\mathrm{Coim}}}

% Boxtimes
%
\DeclareMathSymbol{\boxtimes}{\mathbin}{AMSa}{"02}

%
% Macros for moduli stacks/spaces
%
\def\QCohstack{\mathcal{QC}\!\mathit{oh}}
\def\Cohstack{\mathcal{C}\!\mathit{oh}}
\def\Spacesstack{\mathcal{S}\!\mathit{paces}}
\def\Quotfunctor{\mathrm{Quot}}
\def\Hilbfunctor{\mathrm{Hilb}}
\def\Curvesstack{\mathcal{C}\!\mathit{urves}}
\def\Polarizedstack{\mathcal{P}\!\mathit{olarized}}
\def\Complexesstack{\mathcal{C}\!\mathit{omplexes}}
% \Pic is the operator that assigns to X its picard group, usage \Pic(X)
% \Picardstack_{X/B} denotes the Picard stack of X over B
% \Picardfunctor_{X/B} denotes the Picard functor of X over B
\def\Pic{\mathop{\mathrm{Pic}}\nolimits}
\def\Picardstack{\mathcal{P}\!\mathit{ic}}
\def\Picardfunctor{\mathrm{Pic}}
\def\Deformationcategory{\mathcal{D}\!\mathit{ef}}

\setlength{\emergencystretch}{3em}
\begin{document}
\title[Stacks Project, Tag 07YY]{498 --- Construction of a universal dual obstruction class}
\maketitle
\noindent Copyright (C) 2005--2025 Johan de Jong. Stacks Project authors. Source: \href{https://stacks.math.columbia.edu/tag/07YY}{Tag 07YY}.

\noindent Editorial difficulty note (not author text). Difficulty H2: The proof builds a split extension from the universal polynomial presentation and lifts its obstruction through the naive cotangent complex. It then adjusts the resulting cohomology class by a controlled ambiguity so that its tangent specialization equals the canonical deformation, using a short exact sequence of deformation situations..

\noindent Editorial provenance note (not author text). History: excerpt from revision \texttt{a0444\allowbreak 6e57e\allowbreak c1fbc\allowbreak 252a8\allowbreak 71afc\allowbreak ec775\allowbreak 2fb28\allowbreak 07b14}, artin.tex, lines 4398--4495. Reference presentation was adapted; original-author context and core support blocks were retained or added. The mathematical wording is unchanged. Licensed under GNU FDL 1.2 or later; no invariant sections or cover texts. See ../licenses/COPYING. Context blocks reproduce author definitions and supporting results; they are not additional counted problems.

\begin{situation}
\label{situation-dual}
Let $S$ be a locally Noetherian scheme. Let $\mathcal{X}$ be a category
fibred in groupoids over $(\Sch/S)_{fppf}$. Assume that
$\mathcal{X}$ has (RS*) so that we can speak of the functor $T_x(-)$, see
Lemma \href{https://stacks.math.columbia.edu/tag/07Y9}{lemma properties lift RS star (Tag 07Y9)}.
Let $U = \Spec(A)$ be an affine scheme of finite type over $S$ which maps
into an affine open $\Spec(\Lambda)$. Let $x$ be an object of $\mathcal{X}$
over $U$. Assume we are given
\begin{enumerate}
\item a complex of $A$-modules $K^\bullet$,
\item a transformation of functors
$T_x(-) \to H^1(K^\bullet \otimes_A^\mathbf{L} -)$,
\item for every deformation situation $(x, A' \to A)$ with kernel
$I = \Ker(A' \to A)$ an element
$o_x(A') \in H^2(K^\bullet \otimes_A^\mathbf{L} I)$
\end{enumerate}
satisfying the following (minimal) conditions
\begin{enumerate}
\item[(i)] the transformation
$T_x(-) \to H^1(K^\bullet \otimes_A^\mathbf{L} -)$
is an isomorphism,
\item[(ii)] given a morphism $(x, A'' \to A) \to (x, A' \to A)$ of deformation
situations the element $o_x(A')$ maps to the element $o_x(A'')$
via the map
$H^2(K^\bullet \otimes_A^\mathbf{L} I) \to
H^2(K^\bullet \otimes_A^\mathbf{L} I')$
where $I' = \Ker(A'' \to A)$, and
\item[(iii)] $x$ lifts to an object over $\Spec(A')$ if and only if
$o_x(A') = 0$.
\end{enumerate}
It is possible to incorporate infinitesimal automorphisms as well, but
we refrain from doing so in order to get the sharpest possible result.
\end{situation}

\begin{lemma}
\label{lemma-dual-obstruction}
In Situation \ref{situation-dual}. Assume furthermore that
\begin{enumerate}
\item[(iv)] given a short exact sequence of deformation situations
as in Remark \href{https://stacks.math.columbia.edu/tag/07YE}{remark short exact sequence thickenings (Tag 07YE)} and
a lift $x'_2 \in \text{Lift}(x, A_2')$ then
$o_x(A_3') \in H^2(K^\bullet \otimes_A^\mathbf{L} I_3)$
equals $\partial\theta$ where
$\theta \in H^1(K^\bullet \otimes_A^\mathbf{L} I_1)$
is the element corresponding to $x'_2|_{\Spec(A_1')}$ via
$A_1' = A[I_1]$ and the given map
$T_x(-) \to H^1(K^\bullet \otimes_A^\mathbf{L} -)$.
\end{enumerate}
In this case there exists an element
$\xi \in H^1(K^\bullet \otimes_A^\mathbf{L} \NL_{A/\Lambda})$
such that
\begin{enumerate}
\item for every deformation situation $(x, A' \to A)$ we have
$\xi_{A'} = o_x(A')$, and
\item $\xi_{can}$ matches the canonical element of
Remark \href{https://stacks.math.columbia.edu/tag/07YC}{remark canonical element (Tag 07YC)} via the given transformation
$T_x(-) \to H^1(K^\bullet \otimes_A^\mathbf{L} -)$.
\end{enumerate}
\end{lemma}

\begin{proof}
Choose a $\alpha : \Lambda[x_1, \ldots, x_n] \to A$ with kernel $J$.
Write $P = \Lambda[x_1, \ldots, x_n]$. In the rest of this proof we work with
$$
\NL(\alpha) = (J/J^2 \longrightarrow \bigoplus A \text{d}x_i)
$$
which is permissible by
Algebra, Lemma \href{https://stacks.math.columbia.edu/tag/00S1}{lemma NL homotopy (Tag 00S1)}
and
More on Algebra, Lemma \href{https://stacks.math.columbia.edu/tag/064I}{lemma derived tor homotopy (Tag 064I)}.
Consider the element
$o_x(P/J^2) \in H^2(K^\bullet \otimes_A^\mathbf{L} J/J^2)$ and consider
the quotient
$$
C = (P/J^2 \times \bigoplus A \text{d}x_i)/(J/J^2)
$$
where $J/J^2$ is embedded diagonally. Note that $C \to A$ is a surjection
with kernel $\bigoplus A\text{d}x_i$. Moreover there is a section
$A \to C$ to $C \to A$ given by mapping the class of $f \in P$ to the class
of $(f, \text{d}f)$ in the pushout. For later use, denote $x_C$ the
pullback of $x$ along the corresponding morphism $\Spec(C) \to \Spec(A)$.
Thus we see that $o_x(C) = 0$.
We conclude that $o_x(P/J^2)$ maps to zero in
$H^2(K^\bullet \otimes_A^\mathbf{L} \bigoplus A\text{d}x_i)$.
It follows that there exists some element
$\xi \in H^1(K^\bullet \otimes_A^\mathbf{L} \NL(\alpha))$
mapping to $o_x(P/J^2)$.

\medskip\noindent
Note that for any deformation situation $(x, A' \to A)$ there exists
a $\Lambda$-algebra map $P/J^2 \to A'$ compatible with the augmentations
to $A$. Hence the
element $\xi$ satisfies the first property of the lemma by construction
and property (ii) of Situation \ref{situation-dual}.

\medskip\noindent
Note that our choice of $\xi$ was well defined up to the choice of an
element of $H^1(K^\bullet \otimes_A^\mathbf{L} \bigoplus A\text{d}x_i)$.
We will show that after modifying $\xi$ by an element of the aforementioned
group we can arrange it so that the second assertion of the lemma is true.
Let $C' \subset C$ be the image of $P/J^2$ under the
$\Lambda$-algebra map $P/J^2 \to C$ (inclusion of first factor).
Observe that
$\Ker(C' \to A) = \Im(J/J^2 \to \bigoplus A\text{d}x_i)$.
Set $\overline{C} = A[\Omega_{A/\Lambda}]$. The map
$P/J^2 \times \bigoplus A \text{d}x_i \to \overline{C}$,
$(f, \sum f_i \text{d}x_i) \mapsto (f \bmod J, \sum f_i \text{d}x_i)$
factors through a surjective map $C \to \overline{C}$. Then
$$
(x, \overline{C} \to A) \to (x, C \to A) \to (x, C' \to A)
$$
is a short exact sequence of deformation situations. The
associated splitting $\overline{C} = A[\Omega_{A/\Lambda}]$ (from
Remark \href{https://stacks.math.columbia.edu/tag/07YE}{remark short exact sequence thickenings (Tag 07YE)}) equals the given
splitting above. Moreover, the section $A \to C$ composed with the map
$C \to \overline{C}$
is the map $(1, \text{d}) : A \to A[\Omega_{A/\Lambda}]$ of
Remark \href{https://stacks.math.columbia.edu/tag/07YC}{remark canonical element (Tag 07YC)}.
Thus $x_C$ restricts to the canonical element $x_{can}$ of
$T_x(\Omega_{A/\Lambda}) = \text{Lift}(x, A[\Omega_{A/\Lambda}])$.
By condition (iv) we conclude that $o_x(P/J^2)$ maps to $\partial x_{can}$
in
$$
H^1(K^\bullet \otimes_A^\mathbf{L} \Im(J/J^2 \to \bigoplus A\text{d}x_i))
$$
By construction $\xi$ maps to $o_x(P/J^2)$. It follows that
$x_{can}$ and $\xi_{can}$ map to the same element in the
displayed group which means (by the long exact cohomology sequence)
that they differ by an element of
$H^1(K^\bullet \otimes_A^\mathbf{L} \bigoplus A\text{d}x_i)$
as desired.
\end{proof}
\end{document}
